\section{The weight--clock center tower}\label{sec:center}

For a pure weight-\(w\) motive, the standard completed reflection is
expected to relate \(u\) and \(w+1-u\)
\cite[§1.3]{Serre1970}.  Deligne's conventions place the completed
functional equation and Tate shifts in equations (1.2.3) and (3.1.2),
respectively \cite{Deligne1979}.  We use those conventions only to compute
the \emph{forced center}.  Except at \(n=2\), we do not promote the
expected reflection to a theorem.

\subsection{The canonical index begins at one}

Equation \eqref{eq:cp-denominator} contains the exact identity
\[
 \frac1{p-1}=p^{-1}+p^{-2}+p^{-3}+\cdots.
\]
Thus the canonical index is
\begin{equation}\label{eq:clock}
 j\ge1,\qquad u_{n,j}=ns+j.
\end{equation}
If \(ns'+j=w+1-(ns+j)\), then
\[
 s'=\frac{w+1-2j}{n}-s.
\]
The center of this reflection is
\begin{equation}\label{eq:center-map}
 s_{n,j}(w)=\frac{(w+1)/2-j}{n}.
\end{equation}

For \(j=1\), the three rows are:
\begin{center}
\begin{tabular}{@{}cccc@{}}
\toprule
weight & \(n=2\) & \(n=3\) & \(n=4\)\\
\midrule
\(w=0\)&\(-1/4\)&\(-1/6\)&\(-1/8\)\\
\(w=1\)&\(0\)&\(0\)&\(0\)\\
\bottomrule
\end{tabular}
\end{center}
For \(w=1\) and \(j\ge2\),
\[
 s_{n,j}(1)=-\frac{j-1}{n}.
\]
Therefore the leading odd rail aligns, the leading even rail bifurcates,
and the full tower has no common factorwise standard pure-motive center.

\subsection{An unconditional two-factor witness}

The factorization \eqref{eq:n2-factor} contains two completed factors whose
analytic equations are known.  The Dedekind zeta factor has weight zero and
standard center \(u=1/2\); under \(u=2s+1\), its center is
\[
 s=\frac{1/2-1}{2}=-\frac14.
\]
The curve \(H^1\) factor has weight one and standard center \(u=1\);
under the same clock its center is \(s=0\).  Thus a factorwise mismatch is
already a theorem at \(n=2\), independent of any conjectural \(O_3\) or
\(O_4\) functional equation.

\begin{remark}[Zero is not the obstruction]
The completed Riemann function satisfies
\(\xi(s)=\xi(1-s)\); writing \(\Xi(z)=\xi(1/2+z)\) gives a reflection
about \(z=0\) \cite[§§25.4 and 25.10]{DLMF}.  A zero-centered variable is
therefore entirely compatible with an RH formulation.  The obstruction
here is the simultaneous presence of distinct factorwise centers under the
same locked clock.
\end{remark}

\subsection{Tate and half-shift invariance}

Let \(V(k)\) be an integral Tate twist.  It has weight \(w-2k\).  To
preserve the same local source coefficient, the variable intercept changes
from \(j\) to \(j-k\), because
\(L(V(k),u)=L(V,u+k)\).  Hence
\begin{equation}\label{eq:twist-invariance}
 \frac{(w-2k+1)/2-(j-k)}n
 =\frac{(w+1)/2-j}{n}.
\end{equation}
The center is invariant.

Automorphic unitary normalization and \(C\)-group conventions legitimately
use half-sum shifts \cite{BuzzardGee2010,SST2014}.  Formally setting
\(k=1/2\) in \eqref{eq:twist-invariance} still leaves the center fixed.
If one changes the weight but holds \(u\) fixed, however, the eigenvalue is
multiplied by \(p^{\mp1/2}\).  That operation changes the frozen H\'enon
coefficient and is not a repair of the same object.

\begin{corollary}\label{cor:center-no-go}
No combination of consistent Tate relabelings gives the complete
\(E_n/O_n\) tower a common factorwise standard pure-motive center.
\end{corollary}

The corollary does not exclude a nonfactorwise identity in which additional
source-derived factors cancel or reorganize the separate completions.
